% Generated draft: requires mathematical and source review.
\documentclass[11pt]{article}
\usepackage[margin=25mm]{geometry}
\usepackage{fontspec}
\setmainfont{Latin Modern Roman}
\usepackage{amsmath,amssymb,mathtools,unicode-math}
\setmathfont{Latin Modern Math}
\usepackage{enumitem,needspace,xurl,hyperref}
\setlength{\parindent}{0pt}
\setlength{\parskip}{5pt}
\setlength{\emergencystretch}{3em}
\newcommand{\N}{\mathbb{N}}
\newcommand{\Z}{\mathbb{Z}}
\newcommand{\Q}{\mathbb{Q}}
\newcommand{\R}{\mathbb{R}}
\newcommand{\C}{\mathbb{C}}
\newcommand{\F}{\mathbb{F}}
\newcommand{\A}{\mathbb{A}}
\newcommand{\PP}{\mathbb{P}}
\newcommand{\E}{\mathbb{E}}
\newcommand{\eps}{\varepsilon}
\newcommand{\dd}{\,\mathrm{d}}
\begin{document}
\section*{UnsolvedMath ID: 800109}
\textbf{Source problem code:} catalogue entry
\subsection{Definitions and mathematical statement}
Let \(S^1=\R/\Z\) with its standard orientation, and let \(G=\mathrm{Diff}_+^\infty(S^1)\) be the group, under composition, of orientation-preserving smooth diffeomorphisms of \(S^1\) with smooth inverses. For \(g,h\in G\), define their commutator by
\[
[g,h]=g\circ h\circ g^{-1}\circ h^{-1}.
\]
The question is whether
\[
\forall f\in G\quad \exists g,h\in G\quad f=[g,h].
\]
Thus the requested factors must belong to the same smooth, orientation-preserving circle diffeomorphism group.

For the supplied background, write \(R_\alpha(x)=x+\alpha\) on \(\R/\Z\). The rotation number of \(f\in G\) is
\[
\rho(f)=\left(\lim_{n\to\infty}\frac{F^n(x)-x}{n}\right)\bmod\Z,
\]
where \(F:\R\to\R\) is an increasing lift satisfying \(F(x+1)=F(x)+1\); the resulting class is independent of the lift and of \(x\). Here a rotation number is called Diophantine if a real representative \(\alpha\) satisfies \(\operatorname{dist}(q\alpha,\Z)\geq Cq^{-\tau}\) for every integer \(q\geq1\), for some constants \(C>0\) and \(\tau>0\).

The supplied source attributes to Michel Herman the theorem that every element of \(G\) with Diophantine rotation number is smoothly conjugate to a rotation. It further states that continuity of rotation number yields a factorization of every element of \(G\) as a rotation times an element smoothly conjugate to a rotation, and that every element is consequently a product of two commutators in \(G\). These are supplied background assertions; the question asks whether one commutator always suffices. The source record labels this question open, and this status is not independently verified here.
\subsection{Short English statement}
Does every orientation-preserving smooth diffeomorphism of the circle have the form \(g\circ h\circ g^{-1}\circ h^{-1}\), where \(g\) and \(h\) are also orientation-preserving smooth circle diffeomorphisms? The supplied source states that two commutators suffice and asks whether one always suffices.
\subsection{Sources}
\begin{itemize}
\item Local UnsolvedMath registry: permanent ID \(800109\), exact source code \texttt{catalogue entry}, supplied complete source record. The record attributes the problem to Nariman and Yazdi, with proposed year 2025, and labels it open.
\item Nariman and Yazdi, \emph{Problem List on Foliations and Diffeomorphism Groups} (2025), Question 50. Supplied source URL: \url{https://arxiv.org/abs/2503.18274}. The record reports access on 2026-10-16 and extraction from arXiv LaTeX; these are supplied provenance details, not a claim of independent access or verification.
\item The supplied statement attributes the Diophantine rotation-number conjugacy theorem to Michel Herman. No separate bibliographic reference for that theorem is supplied.
\end{itemize}
\end{document}
